\documentclass[10pt,a4paper]{amsart}
\usepackage[margin=19mm]{geometry}
\usepackage{amsmath,amssymb,mathtools}
\usepackage[scr=rsfs,cal=euler]{mathalfa}
\usepackage{xfrac}
\usepackage[unicode,hidelinks,bookmarks=false]{hyperref}
\newcommand{\ie}{i.e.\ }
\newcommand{\cf}{cf.\ }
\newcommand{\resp}{respectively\ }
\newcommand{\Subsection}{\subsection}
\providecommand{\nobreakdash}{\nobreak\hbox{-}}
\setlength{\emergencystretch}{2em}
\multlinegap0pt
\usepackage{bm}
\usepackage{amssymb}
\usepackage{xcolor}
\usepackage{dsfont}
\usepackage{mathtools}
\usepackage[shortlabels]{enumitem}
\newtheorem{lem}{Lemma}
\usepackage{cleveref}
\crefname{lem}{Lemma}{Lemmas}
\title{Sharp unconditional moment bounds for products of $L$-functions}
\author{Markus Val{\aa}s Hagen}
\date{Source: arXiv:2609.11619v1 (CC BY 4.0)}
\begin{document}
\maketitle

\noindent\emph{Source and licence.} Sharp unconditional moment bounds for products of $L$-functions. Version arXiv:2609.11619v1 (CC BY 4.0), distributed by arXiv under the Creative Commons Attribution 4.0 International licence, http://creativecommons.org/licenses/by/4.0/, as recorded in the arXiv licence metadata for this exact version and shown on the abstract page https://arxiv.org/abs/2609.11619v1. Full licence notice: \texttt{licenses/arxiv-2609.11619-CC-BY-4.0.txt}.\par\smallskip

\noindent\textit{Editorial scope.} Counted unit: the article's lem oneOverNResidualMeanLemma (source0.tex lines 802--819). The article's complete original proof of it is delivered with it (source0.tex lines 1923--2069, one \texttt{proof} environment of 147 lines, which the article places away from the statement). No mathematics is written, repaired or re-derived: the statement and the proof are the article's own lines, in the article's own notation, and no retained line is substituted or re-flowed.  Retained uncounted as the source-local context the proof needs: lines 309--313; lines 782--789; lines 1342--1348; lines 1411--1413; lines 1415--1418.  Nothing was omitted from any retained span, so the package carries no omission record.  No retained line contains a citation, so the wrapper carries no bibliography and no bibliography entry.  No retained line contains a figure environment, an image-inclusion command or drawing code, so no image is delivered and the package declares no asset dependency.  Editorial formatting only, in addition: an AMS \texttt{amsart} preamble that restates the article's own declarations and macros used by the retained lines, namely the environments \texttt{lem} on separate counters, together with the standard packages those lines need.  The retained environments print in this document as Lemma 1 in order of appearance, whereas the article prints the same items with the numbering of its own full document.  The complete native source of the article as distributed in the arXiv e-print for this exact version is bundled unchanged as \texttt{sources/210037/source0.tex}.
		\begin{equation}\label{unweightedVarianceComparison}
			\sum_{T_{j-1}<p\leq T_j}\frac1p
			\leq D_0\left(1+\max_a
			\sum_{T_{j-1}<p\leq T_j}\frac{|\lambda_a(p)|^2}{p}\right).
		\end{equation}

		\begin{equation}\label{residualMajorantParameters}
			\begin{split}
				m_{\rm res}&=\lceil\vartheta_0^{-1}\rceil,\\
				A_{\rm res}&=\sum_{a=1}^r d_a
				\left(c_a+\frac{m_{\rm res}|b_a|}{2}\right),\\
				C_{\rm res}&=2D_0 A_{\rm res}^2+r\log2 .
			\end{split}
		\end{equation}

			\begin{equation}\label{multiblockCoefficientBound}
				\sum_n\frac{|a(n)|^2}{n}
				\ll
				\left(\prod_{\gamma\in\mathfrak F}
				\ell_\gamma!V_\gamma^{\ell_\gamma}\right)
				\exp\left(\sum_{\gamma\in\mathfrak T}V_\gamma\right).
			\end{equation}

		\begin{equation}\label{crossBlockCovariances}
			\sum_{p\leq x}\frac{\lambda_a(p)\overline{\lambda_b(p)}}p=O_{\boldsymbol L}(1).
		\end{equation}

		\begin{equation}\label{primeBlockHierarchy}
			\sum_{\ell>j}P_\ell\ll_{\boldsymbol L,r}\log_{j+1}T
			\leq\varepsilon_*P_j\qquad(2\leq j<\mathscr J).
		\end{equation}

		\begin{lem}\label{oneOverNResidualMeanLemma}
			There is a constant
			$c_0=c_0(\boldsymbol c,\boldsymbol L,n,\mathscr A, \widehat{Y})>0$ such that, for
			every $\boldsymbol\nu\neq
			(\mathscr J+1,\ldots,\mathscr J+1)$,
			\begin{align}
				&\int_T^{2T}
				|\mathcal N_{\boldsymbol\nu}(\tfrac12+it;
				\boldsymbol\beta^{(0)})
				\mathcal Q_{\boldsymbol\nu}(\tfrac12+it)|^2
				R_{\boldsymbol\nu}(t)^{1/\vartheta_0}\,dt\nonumber\\
				&\qquad\ll_{\boldsymbol c,\boldsymbol L,n,\mathscr A, \widehat{Y}}
				T(\log T)^{\kappa(\boldsymbol c)}
				\exp\left(-c_0
				\sum_{\nu_a\leq\mathscr J}K_{\nu_a}\right).
				\label{oneOverNResidualMeanEstimate}
			\end{align}
		\end{lem}

		\begin{proof}[Proof of \Cref{oneOverNResidualMeanLemma}]
			Fix $\boldsymbol\nu\neq
			(\mathscr J+1,\ldots,\mathscr J+1)$ and put
			$\nu_{\min}=\min_a\nu_a$.  Use the constants
			$m_{\rm res},A_{\rm res},C_{\rm res}$ in
			\eqref{residualMajorantParameters}, fixed before $C_{\rm cut}$ and
			$\widehat Y$.  For every $x\geq0$,
			\begin{equation}\label{residualIntegerMajorant}
				x^{2/\vartheta_0}\leq1+x^{2m_{\rm res}}.
			\end{equation}
			Apply this to the exceptional factors, obtaining sums of squared moduli
			of Dirichlet polynomials.
			Before $\nu_{\min}$ all factors belong to the common good prefix.
			The coefficient estimate \eqref{multiblockCoefficientBound}, with
			coefficients $c_a\lambda_a(p)$ and no fixed-power labels, gives
			\begin{equation}\label{residualCommonPrefixNorm}
				\left\|
				\prod_{a=1}^r\prod_{2\leq j<\nu_{\min}}
				\mathcal N_{a,j}(s;2c_a)
				\right\|^{\,2}
				\ll(\log T)^{\kappa(\boldsymbol c)}.
			\end{equation}
			Here the diagonal sum is at most $\kappa(\boldsymbol c)\log_2T+O(1)$,
			and \eqref{crossBlockCovariances} bounds the cross terms after summing
			the common prefix.
			For $j\geq\nu_{\min}$ put
			\[
			\mathcal R_j=\{a:\nu_a\leq j\},\qquad
			\mathcal F_j=\{a:\nu_a=j\},\qquad f_j=|\mathcal F_j|,
			\]
			and let
			\[
			U_j=\sum_{T_{j-1}<p\leq T_j}\frac1p,\qquad
			S_j(s)=\sum_{T_{j-1}<p\leq T_j}p^{-s},
			\]
			with the same unramified-prime convention as before.  By
			\eqref{unweightedVarianceComparison}, $U_j\leq D_0P_j$.
			There are at most $2^{|\mathcal R_j|}\leq2^r$ selections from
			\eqref{residualIntegerMajorant} on this block. The Taylor factors in
			each selected Dirichlet polynomial are
			\[
			\prod_{a:j<\nu_a}\mathcal N_{a,j}(s;2c_a)
			\prod_{a\in E}\mathcal N_{a,j}(s;b_a)^{m_{\rm res}},
			\qquad E\subseteq\mathcal R_j.
			\]
			Since $|\lambda_a(p)|\leq d_a\leq2$, the absolute coefficients of their product
			are dominated coefficientwise by those of
			$\exp(A_{\rm res}S_j(s))$, independently of $K_j$.
			Let $\mathcal A_{j,E}$ be their product with the factors of
			$\mathcal Q_{\boldsymbol\nu}$ on this block. Put $L_j=f_j\ell_j$.
			The product of these latter factors has absolute coefficients bounded by those of
			\[
			\frac{2^{L_j}S_j(s)^{L_j}}
			{(C_{\rm cut}P_j)^{L_j}}.
			\]
			For an integer $L\geq0$ and $A\geq0$, write
			\[
			S_j(s)^L e^{AS_j(s)}=\sum_n\frac{c(n)}{n^s}.
			\]
			If $n$ is supported on this block and $\Omega(n)=L+k$, then
			\[
			c(n)=\frac{(L+k)!A^k}{k!}
			\prod_{p^e\Vert n}\frac1{e!}\qquad(k\geq0),
			\]
			while $c(n)=0$ when $\Omega(n)<L$. The multinomial theorem and
			Leibniz's rule give
			\begin{align}
				\sum_n\frac{|c(n)|^2}{n}
				&\leq\sum_n\frac{|c(n)|^2}{n}\prod_{p^e\Vert n}e!
				\nonumber\\
				&=\sum_{k\geq0}\frac{(L+k)!}{(k!)^2}
				A^{2k}U_j^{L+k}\nonumber\\
				&=e^{A^2U_j}(L!)^2
				\sum_{h=0}^L
				\frac{U_j^{L-h}(A^2U_j^2)^h}{(L-h)!(h!)^2}
				\nonumber\\
				&\leq e^{A^2U_j}L!U_j^L e^{2A\sqrt{LU_j}}.
				\label{residualCoefficientBound}
			\end{align}
			The last inequality follows from $L!/(L-h)!\leq L^h$ and
			$\sum_{h\geq0}y^{2h}/(h!)^2\leq e^{2y}$.
			If $f_j\geq1$, apply \eqref{residualCoefficientBound} with $L=L_j$ and
			$A=A_{\rm res}$, multiply by $4^{L_j}(C_{\rm cut}P_j)^{-2L_j}$,
			and use $U_j\leq D_0P_j$. This bounds $\|\mathcal A_{j,E}\|^2$ by
			\[
			\exp(A_{\rm res}^2D_0P_j)
			\frac{L_j!(4D_0P_j)^{L_j}}
			{(C_{\rm cut}P_j)^{2L_j}}
			\exp\!\left(2A_{\rm res}\sqrt{L_jD_0P_j}\right).
			\]
			Use $L_j!\leq L_j^{L_j}$ and
			$L_j\leq rC_{\rm cut}P_j$, and put
			$\tau=\log(C_{\rm cut}/(4D_0r))\geq16$.
			The logarithm of this bound is at most
			\[
			A_{\rm res}^2D_0P_j-\tau L_j
			+2A_{\rm res}\sqrt{L_jD_0P_j}
			\leq 2A_{\rm res}^2D_0P_j-\frac{\tau}{2}L_j.
			\]
			Here we used
			$2A\sqrt{LU}\leq(\tau/2)L+(2/\tau)A^2U$.
			Since $C_{\rm cut}\geq2$ and $P_j\geq1$,
			$L_j\geq f_jC_{\rm cut}P_j/2$.  Therefore, after summing the
			at most $2^r$ selections,
			\begin{equation}\label{oneOverNLocalBadBlockCoefficientBound}
				\sum_{E\subseteq\mathcal R_j}
				\|\mathcal A_{j,E}\|^{\,2}
				\leq
				\exp\!\left(C_{\rm res}P_j
				-4C_{\rm cut}f_jP_j\right).
			\end{equation}
			For $f_j=0$, use the $L=0$ case of
			\eqref{residualCoefficientBound}.  The bound is independent of
			the truncation orders $K_j$, $\widehat Y$, and $\Lambda$.
			Different blocks have disjoint prime supports, so the sums of squared
			coefficients divided by their indices factor over the blocks.  Every selected polynomial has length
			at most $T^{\varepsilon_0}$ by the length calculation, which
			includes the fixed power $m_{\rm res}$.  Montgomery--Vaughan,
			\eqref{residualCommonPrefixNorm}, and
			\eqref{oneOverNLocalBadBlockCoefficientBound} give
			\begin{align*}
				&\int_T^{2T}
				|\mathcal N_{\boldsymbol\nu}(\tfrac12+it;
				\boldsymbol\beta^{(0)})
				\mathcal Q_{\boldsymbol\nu}(\tfrac12+it)|^2
				R_{\boldsymbol\nu}(t)^{1/\vartheta_0}\,dt\\
				&\qquad\ll T(\log T)^{\kappa(\boldsymbol c)}
				\exp\!\left(
				C_{\rm res}\sum_{j\geq\nu_{\min}}P_j
				-4C_{\rm cut}\sum_{\nu_a\leq\mathscr J}P_{\nu_a}
				\right).
			\end{align*}
			By \eqref{primeBlockHierarchy},
			\[
			\sum_{j\geq\nu_{\min}}P_j
			\leq\sum_{\nu_a\leq\mathscr J}\sum_{j\geq\nu_a}P_j
			\leq2\sum_{\nu_a\leq\mathscr J}P_{\nu_a}.
			\]
			Since $C_{\rm cut}\geq C_{\rm res}$ and
			$K_j\leq11C_{\rm cut}\widehat YP_j$, the last exponent is at most
			\[
			-2C_{\rm cut}\sum_{\nu_a\leq\mathscr J}P_{\nu_a}
			\leq-\frac2{11\widehat Y}\sum_{\nu_a\leq\mathscr J}K_{\nu_a}.
			\]
			This proves \eqref{oneOverNResidualMeanEstimate} with
			$c_0=2/(11\widehat Y)$, independently of $\Lambda$.
		\end{proof}

\end{document}
